\begin{papercard}{stoyanov2015}
\cf{Stages} \stagetag{A} \stagetag{W} \stagetag{I}. Development and first reliability test of the Mobile App Rating Scale (MARS), the most used expert app-quality instrument. A model for building an \emph{anchored} rating instrument from literature criteria and piloting it to reliability.

\cf{Aim} Develop a reliable multidimensional scale to classify and rate mHealth app quality, since star ratings and popularity say little about quality \pc{stoyanov2015}{2}{Selecting apps on the basis of popularity yields}.

\cf{Instrument development} Literature search (2000--2013; PsycINFO, ProQuest, EBSCOhost, IEEE Xplore, Web of Science, ScienceDirect), plus key websites, hand-searching and grey literature, for explicit web/app quality criteria \pc{stoyanov2015}{2}{A comprehensive literature search was conducted to identify}. An expert panel (psychologists, interaction designers, developers, mHealth professionals) classified the criteria into categories and wrote items and descriptors, iterating to full panel agreement \pc{stoyanov2015}{2}{psychologists, interaction and interface designers and}. Of 427 criteria from 25 sources, 56 were duplicates and 22 irrelevant; the remaining 349 were grouped into six categories (Table 1) \pc{stoyanov2015}{3}{A total of 427 criteria were extracted, 56 were}.

\cf{Instrument structure} A classification section (descriptive only: price, platform, ratings, login, password, sharing, target age, health focus), 19 objective items in four subscales (engagement, functionality, aesthetics, information quality), and subjective items \pc{stoyanov2015}{3}{The classification category collected descriptive information}. Each item uses a 5-point scale (Inadequate to Excellent) with \emph{written descriptors for every anchor} and a ``Not applicable'' option where needed. Scores are \emph{means}, not sums, so N/A does not distort them \pc{stoyanov2015}{3}{Mean scores instead of total scores are used because an item}. The form itself shows the anchors, e.g.\ for item 11 (graphics) \pc{stoyanov2015mars}{4}{Graphics: How high is the quality}.

\cf{Testing sample} iTunes search on 19 September 2013 with Boolean terms for mental health and well-being. Inclusion: English, free, Australian store, selected categories \pc{stoyanov2015}{2}{A systematic search of the Apple iTunes store was conducted}. 1{,}533 retrieved; after exclusions by category, title keywords and description, 405 remained, from which 60 were \emph{randomly} selected \pc{stoyanov2015}{4}{A total of 1533 apps were retrieved from the iTunes search}.

\cf{Pilot and revision} Two expert raters each trialled the first 10 apps for at least 10 minutes, rated them independently, then met to resolve ambiguities and revise the scale \pc{stoyanov2015}{3}{trialled each of the first 10 apps for a minimum of 10 minutes}. Pilot ICC was only .57, which led to concrete revisions: N/A removed from engagement items, ``visual information'' clarified, ``young people'' defined as 16--25, the goals item reworded, and an explicit rule on when missing information is N/A vs.\ a flaw \pc{stoyanov2015}{4}{On attempting to rate the initial ten apps}.

\cf{Reliability design} The remaining 50 apps were rated independently. A sample-size argument: at least 41 apps estimate ICC within .15 of .80 with 87\,\% assurance \pc{stoyanov2015}{3}{A minimum sample size}. ICC model: two-way mixed, average measures, absolute agreement, which accounts for proximity rather than exact equality \pc{stoyanov2015}{3}{A two-way mixed effects, average}. Results: total ICC .79, $\alpha$ .90. Subscale ICCs ranged from .50 (functionality) to .80 \pc{stoyanov2015}{4}{Independent ratings on the overall MARS total score of the}. Star ratings correlated only moderately with MARS ($r=.55$, $n=15$ apps with enough ratings) \pc{stoyanov2015}{4}{Only 15 of the 50 mental health and well being apps extracted}.

\cf{Recommendations for raters} Complete training. With multiple raters, agree on the target group, clarify ambiguous items, and decide which items are relevant to the domain. Pilot until reliability or consensus is adequate. Explore content and functions in detail \pc{stoyanov2015}{7}{It is recommended that MARS raters complete a training}.

\cf{Limitations stated} Few peer-reviewed sources. Ease of use and navigation reached only ICC .50. The evidence-base item was untested \pc{stoyanov2015}{7}{items on the functionality subscale (ease of use and navigation)}. Privacy and security were named as future work \pc{stoyanov2015}{7}{evaluate the safety of mHealth apps in terms of the quality}.

\cf{Use for this thesis} A template for the codebook as an instrument: derive items from literature criteria, use expert-panel classification, \emph{write anchor descriptors}, pilot on a separate set, revise with documented changes, then test reliability on a new sample with a justified size. The pilot ICC of .57 and its revisions show why a \emph{revision log} matters. Also a choice of agreement statistic: ICC suits ordinal ratings, while nominal harm codes need $\kappa$/$\alpha$ instead. Quality is not harm communication.
\end{papercard}


\begin{papercard}{stoyanov2015mars}
\cf{What it is} The 7-page English MARS rating form from Stoyanov et al.\ (2015): a classification block, four objective quality sections, a subjective section and an optional app-specific section \pc{stoyanov2015mars}{1}{The Classification section is used to collect descriptive and technical information about the app}. It is an instrument, not a study; the reliability results are in the article card.

\cf{Classification block} Name, ratings and number of ratings for this and all versions, developer, version, last update, costs, platform, brief description, focus, theoretical background or strategies, affiliation (unknown, commercial, government, NGO, university), age group and technical aspects (sharing, community, password protection, login, reminders, web access). The only privacy-related fields are ``allows password-protection'', ``requires login'' and ``allows sharing'' \pc{stoyanov2015mars}{1}{Allows password-protection}.

\cf{Quality sections} ``All items are rated on a 5-point scale'' with a written descriptor for every level \pc{stoyanov2015mars}{2}{All items are rated on a 5-point scale}. Section A engagement has 5 items (entertainment, interest, customisation, interactivity, target group) \pc{stoyanov2015mars}{2}{Entertainment: Is the app fun/entertaining to use}; B functionality 4; C aesthetics 3; D information 7: accuracy of the store description \pc{stoyanov2015mars}{4}{Accuracy of app description (in app store): Does app contain what is described}, goals, quality and quantity of information, visual information, credibility \pc{stoyanov2015mars}{5}{Credibility: Does the app come from a legitimate source}, and evidence base \pc{stoyanov2015mars}{5}{Evidence base: Has the app been trialled/tested}. Items that do not apply are scored N/A and excluded from the section mean \pc{stoyanov2015mars}{5}{Exclude questions rated as}. Section E has four subjective items: recommendation, expected use in 12 months, willingness to pay, star rating \pc{stoyanov2015mars}{6}{What is your overall star rating of the app}. The score sheet reports the four section means, an app quality mean and a separate subjective score. Section F offers six adjustable app-specific items (awareness, knowledge, attitudes, intention to change, help seeking, behaviour change) on 5-point agreement scales \pc{stoyanov2015mars}{7}{These added items can be adjusted and used to assess the perceived impact of the app on the user}.

\cf{Observations on the form} (1) No data-protection or privacy dimension exists among the 19 objective items; the thesis therefore needs privacy instruments (APA, LaMonica, Hutton, Huckvale) beside MARS. (2) The evidence-base descriptors assign level 4 to ``1--2 RCTs'' and level 5 to ``$>$ 3 high quality RCTs'', which leaves exactly three trials unassigned. (3) The credibility item rewards ``nationally competitive government or research funding'' and names Australian bodies, which is context-specific. (4) The item stems are anchored but remain judgement-heavy (``mostly boring'', ``appears correct'').

\cf{Use for this thesis} Copy the \emph{form pattern}: a descriptive classification block, items with an anchor for every level, N/A handling excluded from means, and a separate subjective section so that opinion does not contaminate the objective score. Reuse the classification fields (version, last update, affiliation type, cost, platform) in the app sheet. Do not use MARS scores as a measure of privacy or harm communication.
\end{papercard}
